\subsection{Receiver-centered conditional exchange protocol}
We use two development geometries, a near-contact object with Gaussian material coordinates and a layered object with voxel coordinates. A fixed extension uses a smooth Gaussian object and an asymmetric multi-object voxel target. Each supplies the initial, third-update, and seventeenth-update linearizations, with shared complex current dimensions $k=4,8,16$. The four objects therefore give 36 correlated state--budget cells. All reuse existing noiseless trajectories; states and budgets are not independent objects.

At each state a declared dictionary combines 32 receiver directions, 32 internal-propagation directions, 32 Fourier directions, 32 reduced primal--dual directions, 16 residual-current directions, and ten seeded random currents. Its 154 atoms have a lossless common embedding of rank 128. That embedding caches $L,S,B$ projections; the retained model still has exactly $k$ columns. This is a newly declared dictionary, not the original selector experiment's dictionary. The 1-swap teacher enumerates all feasible replacements in this domain: 600, 1168, and 2208 at the receiver starting spaces for $k=4,8,16$. It is a complete local enumeration, not a global optimum over all current subspaces.

The online proposal uses twelve incoming atoms, balanced across source families without complete-objective labels, and considers each deletable logical atom. It computes endpoint steps and reduced-anchor gains from a separate 64-column model. It then checks at most two finalists by the complete directional formula, accepts a resolved positive gain, and recomputes the conditional neighborhood. At most three actions are accepted. The exact-score teacher, unchecked surrogate, verified proposal, unchanged receiver, and a deletion/addition/swap policy with rank at most $k$ are recorded separately. A fixed restricted 2-swap domain supplies interaction diagnostics. Singular common cores trigger direct endpoint evaluation; endpoint stability determines feasibility.

The material state, full forward residual, injection, prior term $\ell$, damping, and physical coordinates are common to all methods. Full Jacobians and full-reference optimal steps are used only for offline diagnostics. Gaussian spaces have 54 real material coordinates; voxel spaces have 3456. The former allow complete information spectra, while voxel information diagnostics use sixteen fixed common directions. The material scale is $P=10^{-5}I$, separately recorded from LM damping. No measured noise covariance is fitted.

The reported risk is the complete frozen-quadratic gap relative to the paid reference solve. Its reference residual and floating-point floor accompany every cell. Ratios first sum paired risks within each object and only then aggregate objects. The numerical acceptance tolerance is predeclared; it is not a rigorous output-error bound. Selector construction, rejected candidates, failed factorizations, directed RHS, teacher labels and information diagnostics have separate cost records.
